\section{Limitations}
\label{sec:limitations}

% src: FINAL_REPORT.md §9 items 1, 2, 3, 4, 5; §10 item 6 (device groups imposed at feature level); model_zoo.py train_deep (test-based stopping); technical_documentation.md §7.2.1 vs §7.2.6 (single-run variability)
All results come from simulation. Spatial consistency, the property that gives a measurement sequence its information, is a feature of the channel model here; real channels exhibit it, but also add non-stationarity, hardware impairments and calibration drift that the simulations omit. The trajectories are likewise synthetic: they are generated on a grid graph rather than drawn from measured traces, and lack the road geometry and stop-and-go behavior of real movement.

The two device groups of the macro-cell study share one simulated channel, since every UE is simulated with the same omnidirectional array and the absence of an angle is imposed when features are built. The contrast between the groups is therefore a contrast in available measurements rather than in physical hardware. The $85/15$ population mix is an assumption, which is why every result is reported per group rather than for the blend. Each test split of $47$ users contains only $5$ to $10$ single-antenna devices, so their errors carry much wider seed-to-seed uncertainty than those of multi-antenna devices, and any result resting on the nearest range band rests on a few hundred single-antenna samples at most.
% src: diagnostics_2026_10 user_counts (10, 7, 6, 10, 5)

% src: diagnostics_2026_10 meta.models (xgboost, random_forest, knn); mixed_propagation_2026_09/seed_repeats_all_models_summary.json (neural single-antenna gains not significant or negative)
The channel-group test of Section~\ref{sec:results-mechanism} covers $k$-NN, Random Forest and XGBoost only, so the finding that the angle channels carry the gain is established for these three models and inferred for the neural ones. The single-antenna gain that the tree ensembles obtain from RSS history under mixed propagation is unexplained, and the region-crossing account of Section~\ref{subsec:mixed} rests on a single split.

Because every step of a walk advances exactly one grid cell, history depth measures both the number of samples and the distance travelled, and the two cannot be separated in these data; doing so requires simulations with variable step length or variable sampling rate. The neural models chose their stopping epoch on the test users, which flatters their absolute errors and any comparison with the tree and instance-based models, though not the comparisons between depths of one model. Finally, the controlled-grid results are single runs, and repeated runs of the same configuration differ by two to three percentage points, so only their larger effects are interpreted.
